\section{Conclusiones}
\label{sec:conclusiones}

Se construyo y valido experimentalmente un modelo dinamico de emisiones de CO$_2$ en la bioconversion estatica con \textit{Hermetia illucens}, como extension del modelo DEB de Eriksen con tres componentes nuevos: un interruptor de prepupa, un modulo microbiano del lecho y el balance de gas de la camara estatica. La validacion por descomposicion de fuentes, con dos dietas y controles de alimento solo, mostro que la componente larvaria reproduce las tasas netas dentro de la incertidumbre experimental ($R^2 = 0.983$ y RMSE de 5.0 ppm/min en D4; cotas inferiores y punto limpio satisfechos en D1) y que el interruptor de prepupa es indispensable para capturar el colapso de las emisiones al final del ciclo. El ajuste permitio cuantificar, solo con mediciones de gas, la calidad relativa de la dieta de residuos ($f = 0.817$), los tiempos de prepupa por dieta y una biomasa inicial coherente con la postura de huevos, mientras que el modulo microbiano se calibro de forma independiente contra los controles ($k_{\mathrm{ref}} \approx 0.05$ d$^{-1}$, $Y_{CO_2} \approx 0.65$ a $0.69$). Las lecturas de CH$_4$ se mantuvieron como indicador cualitativo. Las extensiones futuras incluyen la validacion directa de biomasa, un estado de crecimiento microbiano y la incorporacion de O$_2$ y N$_2$O.
